\documentclass[11pt]{article}
\usepackage{amsmath,amssymb,amsthm,hyperref}
\newtheorem{theorem}{Theorem}
\newtheorem{lemma}[theorem]{Lemma}
\begin{document}
\title{Symmetric Corrections Do Not Repair Priority-Chain Compression}
\author{}
\date{November 24, 2026}
\maketitle

\paragraph{Scope.}
This note audits a new proposed restricted construction: a sum of
priority-chain heads together with an arbitrary symmetric correction.
It proves that this construction needs exponentially more than the
desired number of heads even for parity with one unused coordinate.
The proof is written only, not Lean-checked. It is not a lower bound
for the original unrestricted attention model, and changes none of
that model's definitions or admissibility conditions.

\section{The exact limiting family of a priority head}

For a permutation $\pi$ of $[n]$, call
\[
 C_\pi(x)=c_0+\sum_{k=1}^n c_k
                              \prod_{\ell=1}^k x_{\pi(\ell)}
                                                        \tag{1}
\]
a chain polynomial. Its coefficients are arbitrary real numbers.
On a Boolean input whose first zero in order $\pi$ occurs at position
$j$, its value is $v_j=\sum_{k=0}^{j-1}c_k$. At the all-one input
its value is $v_{n+1}=\sum_{k=0}^n c_k$.

For $0<\eta<1$, the single affine ratio
\[
 R_{\pi,\eta}(x)=
 \frac{\eta^n v_{n+1}+
       \sum_{j=1}^n\eta^{j-1}v_j(1-x_{\pi(j)})}
      {\eta^n+\sum_{j=1}^n\eta^{j-1}(1-x_{\pi(j)})}
                                                        \tag{2}
\]
converges to $C_\pi$ on every Boolean input as $\eta\downarrow0$.
At a first zero in position $j$, divide numerator and denominator by
$\eta^{j-1}$; all later terms vanish in the limit. The all-one input
has exactly the value $v_{n+1}$ for every $\eta$. Since the cube is
finite, the convergence is uniform. Every denominator in (2) has
strictly negative slopes and minimum $\eta^n>0$ on the whole cube.
It is consequently admissible in the original affine-ratio normal
form. Thus any strict sign representation by a sum of $H$ chain
polynomials transfers to $H$ actual heads by choosing their parameters
sufficiently small.

An arbitrary symmetric correction is a function $Q(|x|)$ on the cube.
Its unique multilinear interpolant has one common coefficient
$\beta_k$ on every monomial of degree $k$. This follows by permuting
coordinates and using uniqueness of multilinear interpolation. Such
a correction is itself available with $O(n)$ heads, for example by
interpolating its $n+1$ values with reciprocals $1/(b_j+|x|)$ at
distinct positive $b_j$. The corresponding square Cauchy matrix is
invertible. The obstruction below permits the symmetric correction
at no cost, so this implementation detail cannot weaken it.

\section{A stronger coefficient-sparsity obstruction}

\begin{theorem}\label{thm:sparse}
Let $n\ge2$, distinguish coordinate $1$, and put
$k=\lfloor(n-2)/2\rfloor$. Suppose a multilinear polynomial
\[
                         P(x)=\sum_{S\subseteq[n]}c_S x^S
\]
satisfies the strict signs
\[
              (-1)^{x_2+\cdots+x_n}P(x)>0
                    \qquad\text{for every }x\in\{0,1\}^n.       \tag{3}
\]
Suppose also that there are real numbers $\beta_0,\ldots,\beta_n$
such that, in each degree, at most $H$ coefficients $c_S$ differ from
$\beta_{|S|}$. Then
\[
                  H\ge\frac12\binom{n-1}{k}.                   \tag{4}
\]
\end{theorem}

\begin{proof}
Write $U=[n]\setminus\{1\}$. For $T\subseteq U$, let
$a_T=(-1)^{|T|}P(0,\mathbf1_T)>0$. Boolean inversion on the face
$x_1=0$ gives, for every $S\subseteq U$,
\[
 c_S=\sum_{T\subseteq S}(-1)^{|S|-|T|}P(0,\mathbf1_T)
       =(-1)^{|S|}\sum_{T\subseteq S}a_T.                       \tag{5}
\]
In particular, $(-1)^{|S|}c_S>0$. If $j\in U\setminus S$, then
\[
 (-1)^{|S|+1}c_{S\cup\{j\}}
       =\sum_{T\subseteq S\cup\{j\}}a_T
       >\sum_{T\subseteq S}a_T
       =(-1)^{|S|}c_S.                                         \tag{6}
\]
The strict inequality uses the positive terms with $j\in T$.
Applying the same inversion argument on the face $x_1=1$ yields
\[
              (-1)^{|S|}(c_S+c_{S\cup\{1\}})>0.               \tag{7}
\]

Call $S$ exceptional when $c_S\ne\beta_{|S|}$. For every
$k$-element $S\subseteq U$ and every $j\in U\setminus S$, at
least one of the three sets
\[
                         S,\quad S\cup\{j\},\quad S\cup\{1\}
                                                                    \tag{8}
\]
is exceptional. Otherwise (5) and (6) say that $\beta_k$ has sign
$(-1)^k$, that $\beta_{k+1}$ has the opposite sign, and that
$|\beta_{k+1}|>|\beta_k|$. Therefore
$\beta_k+\beta_{k+1}$ has sign $(-1)^{k+1}$, contrary to (7).

There are $B d$ pairs $(S,j)$ in (8), where
$B=\binom{n-1}{k}$ and $d=n-1-k$. An exceptional set of degree
$k$ occurs in at most $d$ pairs. An exceptional set of degree
$k+1$ containing coordinate $1$ occurs as $S\cup\{1\}$ in at
most $d$ pairs. One not containing coordinate $1$ occurs as
$S\cup\{j\}$ in exactly $k+1$ possible pairs. Our choice of $k$
ensures $k+1\le d$. Since there are at most $H$ exceptions in each
of these two degrees, all exceptions together cover at most $2Hd$
pairs. Every pair must be covered by the preceding paragraph, so
$Bd\le2Hd$. Since $d>0$, this proves (4).
\end{proof}

\begin{theorem}[Failure of the hybrid proposal]
If an arbitrary symmetric correction plus $H$ chain polynomials
strictly sign-represents parity on coordinates $2,\ldots,n$, then
\[
                  H\ge\frac12\binom{n-1}{\lfloor(n-2)/2\rfloor}
                    =\Omega(2^n/\sqrt n).                     \tag{9}
\]
\end{theorem}

\begin{proof}
Every chain contains at most one monomial in each degree. Thus the
sum of $H$ chain polynomials changes at most $H$ coefficients per
degree from the common coefficient supplied by the symmetric part.
Theorem~\ref{thm:sparse} applies. The final estimate is the usual
central-binomial estimate; the index in (9) is central or one below
central and hence differs only by a bounded factor.
\end{proof}

\paragraph{What this excludes and what it does not.}
The proposed construction used the $n!$ possible priority orders per
head to seek an additional logarithmic saving, with a symmetric term
intended to remove the parity obstruction. Parity on all but one
coordinate defeats that repair. The proof even allows arbitrary
exception locations in each degree, without requiring them to join
into chains. It does not apply to finite nonzero values of $\eta$ in
(2) with cancellation of leading terms, to arbitrary admissible
affine ratios, or to symmetric corrections on arbitrarily chosen
coordinate subsets. Those would be different mechanisms and would
require a new argument.

\end{document}
